%% Apéndice: glosario. Los términos se INTRODUCEN narrativamente en la Parte I;
%% esto es la recapitulación, generada de terminos.json (PLAN-LIBRO.md §2.7).
\chapter{Glosario}
\label{ap:glosario}

Recapitulación de los términos técnicos, en el orden en que el libro los introduce. Cada entrada
lleva la etiqueta de \textbf{de qué lado está lo que nombra}, que es la primera pregunta que hay que
saber contestar de cualquier palabra de este libro.

\begin{center}\footnotesize
\nivelterm{meta} del lado de Lean, donde razonamos\quad
\nivelterm{objeto} del lado del sistema aritmético del que se habla\\[0.35em]
\nivelterm{meta→objeto} una función de Lean que \emph{produce} sintaxis de ese sistema\quad
\nivelterm{ambos} nombra la distinción misma\\[0.35em]
\nivelterm{proyecto} una noción sobre el proyecto, no sobre las teorías
\end{center}

\input{extraido/_glosario.tex}
